\hypertarget{conclusion}{\section{Conclusion}}
\label{sec:conclusion}

Language models have transformed AI, ushering in the paradigm of foundation models.
The reach of modern language models extends well beyond research, with language models being rapidly productionized into consequential and ubiquitous language technologies, which we expect to only increase in the near-term future. 
We lack transparency on language models at present, which is especially concerning given their rapid growth and burgeoning impact: as a community, we do not understand language models in their totality.
For this reason, we have pushed for holistic evaluation in this effort, as we believe holistic evaluation is a critical means for providing the necessary transparency for language models.

Transparency begets trust and standards. 
Viewing benchmarks as models for social change, given they orient the development of AI systems, our broader objective is to transform foundation models from immature emerging technologies to reliable tools that support human flourishing. 
With this objective in mind, we recognize the history and trajectory of AI benchmarking aligns with institutional privilege \citep{koch2021reduced}. 
Benchmarks set the agenda and orient progress: we should aspire for holistic, pluralistic and democratic benchmarks \citep{birhane2022values}. 
Given the understated but significant power of benchmarks to drive change, which in turn indicates benchmark design confers power, we foreground our objectives for \benchmarkname along with its limitations. 
We hope the community will interrogate, adopt, and improve \benchmarkname going forward to actualize the ambition of holistic evaluation.
In this way, we hope holistic evaluations for language models and for other classes of foundation models will give rise to useful, responsible, and societally beneficial technology. 